% Directly authored companion to gamma_access_collaborator_note.md.
\documentclass[11pt]{article}
\usepackage[T1]{fontenc}
\usepackage{amsmath,amssymb,amsthm,newtxtext,newtxmath}
\usepackage[letterpaper,margin=0.76in]{geometry}
\usepackage{graphicx,microtype,xcolor,caption,fancyhdr}
\usepackage[colorlinks=true,linkcolor=blue!40!black,urlcolor=blue!50!black]{hyperref}
\definecolor{ink}{HTML}{263442}
\DeclareMathOperator{\sech}{sech}
\DeclareMathOperator{\sign}{sign}
\DeclareMathOperator{\range}{range}
\theoremstyle{definition}
\newtheorem{theorem}{Theorem}
\newcommand{\repo}{https://github.com/samlayton99/precision-mlps/blob/experiment/frozen-gamma-probe}
\newcommand{\evidence}{\repo/results/checkpoint_D_optimizers/expD36_frozen_gamma_probe/full_sweep/refinements/gamma_access_note}
\graphicspath{{results/checkpoint_D_optimizers/expD36_frozen_gamma_probe/full_sweep/refinements/gamma_access_note/}}
\hypersetup{pdftitle={How gamma controls access to representable target components},
pdfsubject={Finite tanh readout theorem, positive slow target energy, proof, and measured optimization delays}}
\setlength{\parindent}{0pt}
\setlength{\parskip}{4pt plus 1pt}
\setlength{\emergencystretch}{2em}
\setlength{\headheight}{14pt}
\captionsetup{font=small,labelfont=bf,justification=raggedright,singlelinecheck=false}
\pagestyle{fancy}\fancyhf{}
\fancyhead[L]{\small\color{ink}Gamma and optimization access}
\fancyhead[R]{\small Collaborator note}
\fancyfoot[C]{\small\thepage}
\renewcommand{\headrulewidth}{0.3pt}

\begin{document}
\thispagestyle{plain}
{\color{ink}\fontsize{22}{25}\selectfont\bfseries
How gamma controls access to\\representable target components\par}
\vspace{7pt}
\textbf{The proposed paper claim.} A target component can belong to a frozen
tanh dictionary's span and still require millions of readout updates to learn.
A common small slope $\gamma$ suppresses frequency content in the features.
On a finite interval, sampling and boundaries mix those frequencies; retaining
this mixing carries the explicit attenuation into accurate eigenvalues and
target projections. We can therefore identify \textbf{representable but slow
target components}, and predict the resulting delay before running GD.

For a paper with room for one theorem, we recommend the \textbf{finite-interval
tanh theorem} below. It concerns the features actually trained. The periodic
calculation is a motivating example and a step in the proof, rather than a
second headline theorem. Evaluating the result includes a finite spectral
calculation. No decay rate is fitted to the observed training curve.

The figure follows three steps: gamma filters frequency content; the finite
kernel places target energy in positive but slow modes; that spectrum predicts
the observed training trajectory. At gamma 8, the calculation places at least
\textbf{4.3322\% of target energy in representable slow modes}, implying a
\textbf{3.036-million-update necessary time} from that band alone. The complete
spectrum gives a narrow interval containing the measured
\textbf{15,798,313 updates}. Spectral numbers use a disclosed floating-point
allowance; the four primary full-spectrum timing intervals also contain
independently certified archived intervals.

\textbf{Notation.} $J_\gamma$ is the sample-normalized feature matrix;
$K_\gamma=J_\gamma J_\gamma^*$ is its readout kernel. $\mu_j$ is an eigenvalue,
and $\eta\mu_j$ is its per-update learning rate. $E_\gamma(n)$ is relative
residual norm: 1\% means $E_\gamma(n)\le0.01$. Physical slope is $\gamma$;
dimensionless bandwidth is $\lambda=\gamma h$, never an eigenvalue. A hat
denotes the constructed approximation; $\delta_\gamma$ bounds its kernel
error. The band-energy fraction is $p_\gamma$; the periodic reference feature
is $\phi_\gamma$.

\section{Why the periodic example makes the mechanism transparent}
A tanh is exactly a smoothed step:
\[
\tanh(\gamma x)=(\rho_\gamma*\sign)(x),\qquad
\rho_\gamma(t)=\frac\gamma2\sech^2(\gamma t).
\]
The smoothing density has Fourier multiplier
\[
M_\gamma(\omega)=\frac z{\sinh z},\qquad
z=\frac{\pi|\omega|}{2\gamma},\qquad M_\gamma(0)=1.
\]
On a circle, with centers and input measure uniform over the whole period,
translations make the kernel diagonal in Fourier waves. Smoothing multiplies
each feature coefficient by $M_\gamma$, and the kernel pairs two coefficients:
\[
\mu_\omega(\gamma)=\mu_\omega^{\rm step}M_\gamma(\omega)^2.
\]
At high frequency relative to gamma,
\[
M_\gamma(\omega)^2\sim
\left(\frac{\pi|\omega|}{\gamma}\right)^2e^{-\pi|\omega|/\gamma}.
\]
At a fixed step size, this produces an exponential optimization penalty.
Every frequency with
$\mu_\omega^{\rm step}>0$ remains accessible at every $\gamma>0$, but its
learning rate can become extremely small. Frequencies absent from the
periodic step remain inaccessible; for the standard square wave these
include nonconstant even harmonics. Bias supplies the constant mode. The
claim concerns accessible frequencies, not all periodic functions.

Ordinary tanh on an interval is not this periodic population model.
Boundaries and finite sampling can rotate learning directions. The finite
theorem retains these effects instead of treating periodic eigenvalues as
an approximation to the answer.

\section{The one theorem to present}
Fix an interval of length $L>0$, an integer $N\ge3$, and equally spaced core centers $c_j=jh$,
$j=0,\ldots,N-1$, $h=L/N$, and inputs $x_i=ih/q$, where $q$ is a positive
integer and $i=0,\ldots,qN$. Thus $m=qN+1$ includes both endpoints.
Any extra centers stay fixed and are retained explicitly; $W$ counts all
tanh features. Include a bias column equal to $1/\sqrt m$ and define
\[
J_\gamma[i,j]=\frac{\tanh(\gamma(x_i-c_j))}{\sqrt m},\qquad
y_i=\frac{y_i^{\rm phys}}{\sqrt m}.
\]
All tanh features have a common frozen slope $\gamma>0$.
GD minimizes $\frac12\|J_\gamma w-y\|^2$ in the original Euclidean
readout coefficients, starts at $w=0$, and uses $\eta>0$.
The nonzero target $y$ and all input/center locations stay fixed as gamma
changes.

Its relative residual is
$E_\gamma(n)=\|(I-\eta K_\gamma)^ny\|/\|y\|$.
Write $\mu_j$ for eigenvalues in decreasing order, including zeros, and
$u_j$ for an orthonormal eigenbasis. Define cumulative target energy
\[
S_\gamma(t)=\sum_{\eta\mu_j\le t}
\frac{|\langle y,u_j\rangle|^2}{\|y\|^2}.
\]

First synthesize a periodic reference from Fourier coefficients multiplied
by $M_\gamma$. Then retain the gamma-dependent boundary terms, exact boundary
columns, extra centers, bias, and endpoint row. Call the resulting features
$\widehat J_\gamma$ and kernel $\widehat K_\gamma$. Appendix A gives the
scalar formulas and error allowance; no training observation is used.

\begin{theorem}[Gamma-dependent access to representable finite-tanh components]
Under the stated geometry and raw-readout parameterization, the Fourier and
boundary construction gives an explicit allowance $\delta_\gamma$ such that
\[
\|K_\gamma-\widehat K_\gamma\|\le\delta_\gamma,\qquad
|\mu_j(\gamma)-\widehat\mu_j(\gamma)|\le\delta_\gamma.
\]
All gamma dependence is explicit in the frequency multipliers, boundary
factors, and allowance. The analytic allowance decreases geometrically to
zero as the retained boundary expansion increases.

Choose $\eta(\|\widehat K_\gamma\|+\delta_\gamma)\le1$. Let
$\underline S_\gamma(t)$ and $\overline S_\gamma(t)$ be the bounds defined
below on the fraction of target energy in per-update rates at most $t$.
For $0<a<b<1$, set
\[
p_\gamma=\left[\underline S_\gamma(b)-\overline S_\gamma(a)\right]_+.
\]
At least a fraction $p_\gamma$ of target energy lies in the positive-rate
band $(a,b]$. This component is exactly representable on the sampled grid.
Nevertheless, it implies
\[
\boxed{E_\gamma(n)\ge\sqrt{p_\gamma}(1-b)^n.}
\]
\end{theorem}
The theorem identifies a property of the specified target and dictionary.
It does not assume target alignment with a chosen Fourier wave. Gamma changes
both eigenvalues and finite eigenvectors; the spectral calculation retains
both changes when computing target energy.

For explicit energy bounds, compute the eigenvectors $\widehat u_j$ of
$\widehat K_\gamma$ and set
\[
\widehat S_\gamma(t)=\sum_{\eta\widehat\mu_j\le t}
\frac{|\langle y,\widehat u_j\rangle|^2}{\|y\|^2}.
\]
Include the full nullspace. For a chosen buffer $g>0$, valid bounds are
\begin{align*}
\underline S_\gamma(t)
&=\left(\sqrt{\widehat S_\gamma(t-g)}-
\frac{\eta\delta_\gamma}{g}\right)_+^2,\\
\overline S_\gamma(t)
&=\min\!\left\{1,\left(\sqrt{\widehat S_\gamma(t+g)}+
\frac{\eta\delta_\gamma}{g}\right)^2\right\}.
\end{align*}
Buffers may differ between cutoffs. Taking the largest lower and smallest
upper bound over any declared collection remains valid. Subtracting upper
mass at a strictly positive $a$ is essential: it excludes the possibility
that the reported slow energy is merely an unlearnable nullspace component.

The full target need not be exactly representable. The component isolated
by $(a,b]$ is. This precisely separates optimization access from capacity.

\section{Where gamma enters, and what must still be calculated}
The quantitative chain is
\[
\begin{aligned}
\gamma
&\longrightarrow\underbrace{M_\gamma(\omega)\ \text{and boundary factors}}
_{\text{explicit feature construction}}\\
&\longrightarrow\underbrace{\widehat\mu_j(\gamma),\widehat u_j(\gamma)}
_{\text{finite spectral calculation}}
\longrightarrow\underbrace{p_\gamma\ \text{and }E_\gamma(n)}
_{\text{target access and delay}}.
\end{aligned}
\]
This goes beyond observing that gamma changes a matrix. We identify its
feature-level action analytically, retain the geometry that mixes frequencies,
and bound the resulting spectral error. Finite eigenvectors are computed,
not assumed to remain Fourier waves. The boundary correction also depends
on gamma; it is neither a free fit nor an assumed small perturbation.

Two errors must not be confused. Small gamma can make true learning rates
small. Separately, truncation introduces approximation error $\delta_\gamma$,
reducible by retaining more terms. The latter controls prediction precision;
it is not the proposed cause of delay.

The periodic example gives a scalar cap-to-rate inequality immediately.
On the finite interval, the theorem gives quantitative predictions after
solving the explicit corrected spectral problem. It does not assert that
every target or finite eigenvalue improves monotonically with gamma. Such
ordering is unnecessary to explain the measured, target-dependent delay.

\section{Predictions and the three-panel figure}
\begin{figure}[htbp]
\centering
\includegraphics[width=\linewidth]{gamma_access_three_panel.pdf}
\caption{\textbf{Gamma changes optimization access to a fixed target.}
A shows exact analytic periodic attenuation, not measured training data.
Faint guides mark the target's physical frequencies, without identifying
them with finite-kernel eigenvectors.
B shows finite ordinary-tanh target energy in positive rates above
$a=10^{-9}$ and below the horizontal-axis cutoff. Solid curves come from
the original tanh spectrum; shaded bounds come from the Fourier-plus-boundary
construction. The marked cutoff is $b=10^{-6}$. C compares predictions
from the complete constructed spectrum with saved GD checkpoints.
The horizontal 1\% threshold and crossing markers show acquisition time.
Circles are saved GD measurements, diamonds are executed first hits,
and short horizontal segments at 1\% are necessary--sufficient timing
intervals. Colors identify the same four slopes throughout. Spectral bands
use the disclosed FP64 sensitivity allowance, not independent interval
certification. The primary full-spectrum timing intervals contain
previously certified intervals.}
\end{figure}

The primary target is
\[
f(x)=\sin(2\pi x)+\tfrac12\sin(6\pi x)+\tfrac14\sin(10\pi x),
\qquad x\in[-1,1].
\]
The archive uses $L=2$, $N=512$, $q=16$, and $m=8193$. Its 512 core
centers, terminal center, and 23 halo centers on each side give 559 tanh
features plus bias. Common frozen slopes are 8, 12, 16, and 64. Readouts
start at zero; each saved step is curvature-normalized to approximately
$\eta\|K_\gamma\|=0.5$. Comparisons therefore concern those actual normalized
steps, not one common raw step. Translating $[-1,1]$ to $[0,L]$ changes
neither feature differences nor target samples.

At $\gamma=8$, $a=10^{-9}$, $b=10^{-6}$, numerical evaluation gives
\[
\underline S_8(b)=0.0433317719,\qquad
\overline S_8(a)=0.00000929944,\qquad
p_8\ge0.0433224724.
\]
Thus at least 4.3322\% of target energy lies in strictly positive but slow
modes. This band requires at least \textbf{3,035,627 updates} to reach 1\%
relative residual. At gamma 64, even the upper bound on all energy with
rates at most $10^{-6}$ is only about \textbf{0.000162\%}.
These compare the same target, retaining each gamma's own finite directions.

The single-band bound discards variation within that band. Keeping every
rate and target weight gives the necessary--sufficient intervals
\[
\begin{gathered}
\left[15{,}767{,}081,15{,}829{,}763\right],\qquad
\left[186{,}056,186{,}059\right],\\
\left[61{,}792,61{,}793\right],\qquad
\left[16{,}013,16{,}013\right].
\end{gathered}
\]
Measured crossings are \textbf{15,798,313}, \textbf{186,057},
\textbf{61,792}, and \textbf{16,013}. The gamma-8 interval width is
0.397\% of its observed time. It is useful to show both the positive slow
component explaining the obstruction and the full forecast predicting
its duration.

This retrospectively analyzes existing optimization experiments. The
representative target and rate band were chosen during analysis, not
preregistered. Truncations use analytic error bounds; the buffer scan uses
the constructed spectrum without consulting the original target mass.
There is no new training or fitted decay parameter, and no held-out
generalization claim. The wider package has three grid sizes, four slopes,
and five targets, with all 60 reference crossings and all 19 available
executed crossings enclosed.

The spectral bands add a heuristic FP64 arithmetic allowance to the proved
analytic remainder; they are not independent interval certificates.
The four primary full-spectrum timing intervals also contain independently
certified Arb intervals, validating the wider timing intervals by monotonicity.
This inherited validation does not certify the new positive-band energy
numbers. The \href{\repo/docs/gamma_uniform_grid_theorem.md}{full technical
note} records the arithmetic allowance and its tenfold sensitivity check.

\appendix
\section{Explicit construction and proof}
\textbf{The reference and its boundaries.}
Let $s$ be the $2L$-periodic square wave equal to $\sign(d)$ on $(-L,L)$,
and $\phi_\gamma=\rho_\gamma*s$. It is antiperiodic: shifting by $L$ flips
its sign. Its Fourier expansion is
\[
\phi_\gamma(d)=\sum_{k\in\mathbb Z}\frac2{i\pi(2k+1)}
M_\gamma\!\left(\frac{(2k+1)\pi}{L}\right)e^{i(2k+1)\pi d/L}.
\]
Positive gamma gives absolute convergence. Sampling retains aliases:
coefficients of frequencies indistinguishable on the finite grid add
before forming the kernel. Do not substitute population eigenvalues
for this sampled calculation.

The exact image identity is
\[
\phi_\gamma(d)=\tanh(\gamma d)+\sum_{\ell\ge1}(-1)^\ell
\big[\tanh(\gamma(d-\ell L))+\tanh(\gamma(d+\ell L))\big].
\]
Both sides have derivative
$\sum_{\ell\in\mathbb Z}(-1)^\ell\gamma\sech^2(\gamma(d-\ell L))$.
Both vanish at zero and their paired tails converge locally uniformly,
proving the identity. Expanding each image's positive-distance tanh as
a geometric series gives, for $|d|<L$,
\[
\tanh(\gamma d)=\phi_\gamma(d)+
2\sum_{r\ge1}\frac{(-1)^{r-1}}{1+e^{-2r\gamma L}}
\big[e^{-2r\gamma(L-d)}-e^{-2r\gamma(L+d)}\big].
\]
This is explicit finite-interval gamma dependence: the Fourier attenuation
and every boundary coefficient are known functions of gamma. With $d=x-c$,
each exponential separates into a factor of $x$ and one of $c$.
Retaining $p$ terms gives correction rank at most $2p$ on interior columns.

\textbf{A controllable remainder.}
Choose an integer $B$ with $1\le B<N/2$ and a nonnegative integer $p$.
Treat the first and last $B$ core columns exactly.
On the other core columns and all but the endpoint input row, use
$\phi_\gamma$ plus the first $p$ boundary terms. Treat extra centers, bias,
and the endpoint row exactly. This defines $\widehat J_\gamma$ in the
original feature coordinates.

Interior distances have $|d|\le L-Bh$. The finite geometric remainder of
each tanh expansion is at most twice its first omitted exponential.
Summing both sides and all images gives the unnormalized entry-error bound
\[
\epsilon_{B,p}(\gamma)=
\frac{4e^{-2(p+1)\gamma Bh}}{1-e^{-2(p+1)\gamma L}}.
\]
At most $N-2B$ columns have error. Frobenius norm bounds operator norm, so
\[
\epsilon_J=\sqrt{N-2B}\,\epsilon_{B,p}(\gamma),\qquad
\|J_\gamma-\widehat J_\gamma\|\le\epsilon_J.
\]
Every true unnormalized entry has magnitude at most one, including bias.
Thus $\|J_\gamma\|\le\sqrt{W+1}$ and the expanded Gram difference gives
\[
\boxed{\delta_\gamma=2\sqrt{W+1}\,\epsilon_J+\epsilon_J^2.}
\]
This proves the kernel bound; the min--max characterization proves the
ordered eigenvalue bound. If the reference Fourier series is truncated
with uniform entry error $\tau$, replace $\epsilon_{B,p}$ by
$\epsilon_{B,p}+\tau$ before forming $\epsilon_J$ and $\delta_\gamma$.
Floating-point evaluation needs a separate allowance.

The reference's uniform-grid blocks are explicitly diagonalizable by
half-odd Fourier vectors. Restoring boundary columns, the $p$ separable
terms, $W-N$ extra centers, bias, and endpoint row changes the feature
matrix by rank at most $2B+2p+(W-N)+2$. The kernel correction rank is
at most twice this. It is signed and need not be small in norm. The
\href{\repo/docs/gamma_uniform_grid_theorem.md}{full construction proof}
supplies exact alias sums and signed spectral formulas. The numerical full
forecast uses the SVD of the constructed feature matrix; it does not claim
to obtain all target projections from a small eigenvalue-count calculation.

\textbf{Transferring target energy.}
Put $A=\eta K_\gamma$, $\widehat A=\eta\widehat K_\gamma$, and
$d=\eta\delta_\gamma$. Let $P$ project onto eigenvalues of $A$ at most $t$,
and $Q$ onto eigenvalues of $\widehat A$ at most $t-g$. The cross-projection
$X=(I-P)Q$ solves
\[
A_{>t}X-X\widehat A_{\le t-g}=(I-P)(A-\widehat A)Q.
\]
The restricted spectra are separated by at least $g$. The exponential
integral solution of this Sylvester equation gives $\|X\|\le d/g$.
Therefore $\|Qy\|\le\|Py\|+(d/g)\|y\|$, yielding the lower cumulative
bound. Use the projector of $\widehat A$ below $t+g$ and reverse the roles
for the upper bound. This controls subspaces without requiring a gap
between every pair of eigenvalues.

True energy in $(a,b]$ is $S_\gamma(b)-S_\gamma(a)$ and is at least
$[\underline S_\gamma(b)-\overline S_\gamma(a)]_+$.
Every eigenvalue in this band is positive, so its eigenvector belongs to
$\range(K_\gamma)=\range(J_\gamma)$. This proves exact representability
of the isolated component on the input grid.

\textbf{Transferring to training.}
Zero initialization gives $r_n=(I-\eta K_\gamma)^ny$. Under the step
condition, each rate in $(a,b]$ retains at least a factor $1-b$ of residual
amplitude per update.
Square and sum the target coefficients to obtain
$E_\gamma(n)^2\ge p_\gamma(1-b)^{2n}$.
For $p_\gamma>\varepsilon^2$, reaching $E_\gamma(n)\le\varepsilon$ requires
\[
n\ge\left\lceil
\frac{\log(\sqrt{p_\gamma}/\varepsilon)}{-\log(1-b)}
\right\rceil.
\]
The complete constructed spectrum gives
\begin{align*}
\widehat E_\gamma(n)^2
&=\sum_j\frac{|\langle y,\widehat u_j\rangle|^2}{\|y\|^2}
(1-\eta\widehat\mu_j)^{2n},\\
|E_\gamma(n)-\widehat E_\gamma(n)|&\le n\eta\delta_\gamma.
\end{align*}
The sum includes nullspace energy. Telescope the two matrix powers:
each of $n$ terms contains one kernel difference and otherwise contractions.
This proves the last inequality. Thus
$\widehat E_\gamma(n)+n\eta\delta_\gamma\le\varepsilon$ proves acquisition
by $n$, while $\widehat E_\gamma(n)-n\eta\delta_\gamma>\varepsilon$ excludes
it by $n$. This is the origin of the necessary--sufficient timing intervals.

\section{Paper placement and reproducibility}
Use the periodic multiplier as a motivating example, the finite-tanh access
result as the sole numbered theorem, and the three-panel figure as its
empirical evaluation. Keep the scalar construction, transfer proof, full
timing formula, and numerical sensitivity in the appendix. Present common
frozen slope, uniform core centers and samples, fixed target, and raw readout
GD alongside the theorem. Representability means representability on the
training grid.

At grid-relative frequency $\theta=\omega h$, attenuation depends on
$\gamma h$ through $z=\pi|\theta|/(2\gamma h)$. Holding $\gamma h$ fixed
preserves this part of spectral access under refinement. It does not alone
remove raw-kernel width scaling, boundary effects, or step dependence.

The \href{\evidence/gamma_access_data.json}{figure data} records source
hashes, positive-band bounds, GD checkpoints, and the distinction between
analytic and measured quantities. Reproduce the figure with
\begin{quote}\small\ttfamily
python -m experiments.expD36\_frozen\_gamma\_probe.gamma\_access\_figure
\end{quote}
Compile this note with
\begin{quote}\small\ttfamily
latexmk -pdf -outdir=/tmp/gamma-access-latex\\
\hspace*{1em}docs/gamma\_access\_collaborator\_note.tex
\end{quote}
Earlier proof notes and archives remain unchanged. No new GPU training is used.
\end{document}
